\section{Trabalhos relacionados}

Este capítulo será desenvolvido a partir da revisão já iniciada no pré-projeto e no artigo-base. Sua função é posicionar a contribuição da dissertação diante dos estudos mais próximos, comparando explicitamente corpus, idioma, jurisdição, arquitetura, técnicas de avaliação, participação humana e disponibilidade dos artefatos.

\subsection{Sistemas RAG nos domínios jurídico, legislativo e parlamentar}

% TODO: transformar a síntese existente no referencial teórico em comparação
% crítica apoiada pela tabela de literatura do artigo-base.

\subsection{Protocolos, métricas e benchmarks de avaliação de RAG}

% TODO: comparar RAGAS, ARES, LegalBench-RAG, RAGEval e estudos de LLM-as-a-judge.

\subsection{Síntese comparativa e posicionamento da dissertação}

% TODO: inserir quadro comparativo por corpus, língua, pipeline, recuperação,
% geração, avaliação humana, juiz automatizado e abertura dos artefatos.
